\documentclass[a4paper, serif, 10pt]{moderncv}

%% moderncv
% style and color
\moderncvstyle{banking}
\moderncvcolor{black}

%% page layout/margins
\usepackage[top=2.5cm, bottom=2.5cm, left=1.5cm, right=1.5cm]{geometry}

\nopagenumbers{}

%% Babel config for English language
\usepackage[english]{babel}

%% fontspec
\usepackage{fontspec}

\IfFontExistsTF{Linux Libertine}{
  \setmainfont[Ligatures={TeX, Common}, Numbers=OldStyle]{Linux Libertine}
}{}
\IfFontExistsTF{Linux Libertine O}{
  \setmainfont[Ligatures={TeX, Common}, Numbers=OldStyle]{Linux Libertine O}
}{}

%% CTeX
% CJK support, used to show CJK characters in the resume
%
% - fontset=none: disable builtin fontset but instead set the CJK font manually
% - heading=false: disable ctex heading
% - punct=kaiming: use kaiming punctuations styles for CJK
% - scheme=plain: use plain scheme, do not override \normalsize font size
% - space=auto: space settings for CJK characters
%
% ref:
% - http://ctan.mirrorcatalogs.com/language/chinese/ctex/ctex.pdf
\usepackage[UTF8, heading=false, punct=kaiming, scheme=plain, space=auto]{ctex}

\IfFontExistsTF{Noto Serif CJK SC}{
  \setCJKmainfont{Noto Serif CJK SC}
}{}
\IfFontExistsTF{Noto Sans CJK SC}{
  \setCJKsansfont{Noto Sans CJK SC}
}{}

%% line spacing
\usepackage{setspace}
\setstretch{1.125}

%% URL styling - use normal text instead of monospace
\urlstyle{same}

% Auto-underline all links
% Defer until \begin{document} so hyperref is already loaded
\AtBeginDocument{%
  \let\oldhref\href
  \renewcommand{\href}[2]{\oldhref{#1}{\underline{#2}}}%
}

%% Patch \httplink and \httpslink to handle full URLs
% The original moderncv macros always prepend http:// or https:// to URLs.
% This causes issues when the URL already has a protocol (e.g., https://example.com)
% resulting in malformed URLs like https://https://example.com.
% This patch checks if the URL already contains "://" and uses it directly if so.
% Note: We use \str_set:Nx (x-expansion) to expand macros like \@homepage first.
\ExplSyntaxOn
\str_new:N \g_yamlresume_url_str

\RenewDocumentCommand{\httpslink}{O{}m}{
  \str_set:Nx \g_yamlresume_url_str {#2}
  \str_if_in:NnTF \g_yamlresume_url_str {://}
    { \url{#2} }
    { \url{https://#2} }
}

\RenewDocumentCommand{\httplink}{O{}m}{
  \str_set:Nx \g_yamlresume_url_str {#2}
  \str_if_in:NnTF \g_yamlresume_url_str {://}
    { \url{#2} }
    { \url{http://#2} }
}
\ExplSyntaxOff

\name{林彥廷}{}
\title{資深 DevOps 工程師 | 雲端平台與自動化維運專家}
\phone[mobile]{+886 912 345 678}
\email{yenting.lin@example.com}
\homepage{https://www.yentinglin.dev}

\address{台北市大安區和平東路二段 100 號 8 樓, 台北市, 台北市, Taiwan, 106}{}{}

\extrainfo{{\small \faGithub}\ \href{https://github.com/yentinglin}{@yentinglin} {} {} {} • {} {} {} 
{\small \faLinkedin}\ \href{https://www.linkedin.com/in/yentinglin}{@yentinglin} {} {} {} • {} {} {} 
{\small \faStackOverflow}\ \href{https://stackoverflow.com/users/1234567/yenting-lin}{@yenting-lin}}

\begin{document}

\maketitle

\section{Basics}

\cvline{}{資深 DevOps 工程師，擁有 8 年以上雲端平台、自動化維運與可靠性工程經驗。專精於 Kubernetes、Terraform、CI/CD 與可觀測性建置，並擅長推動 DevSecOps 文化與跨團隊協作。
%
\begin{itemize}
\item 主導多叢集 Kubernetes 平台與 GitOps 流程，支援每日數百次部署，將部署成功率提升至 99.9\%
\item 以基礎架構即程式碼（IaC）管理 AWS、GCP 與混合雲環境，降低 40\% 人為操作失誤
\item 建立 SRE 監控與告警體系，將平均修復時間（MTTR）縮短 55\%
\item 具備良好的溝通與文件撰寫能力，樂於分享知識並培育初階工程師
\end{itemize}}
\section{Education}

\cventry{Sep 2014–Jun 2016}
        {Master, 資訊工程學系, Score: 4.1}
        {國立臺灣大學}
        {\url{https://www.ntu.edu.tw/}}
        {}
        {\begin{itemize}
\item 專注於分散式系統與雲端運算，論文主題為「基於容器編排之自動擴展機制」
\item 擔任計算機網路助教，協助課程設計與學生實作指導
\item 參與開源社群，累積 Kubernetes 與自動化維運實務經驗
\end{itemize}
\textbf{Courses}: 分散式系統, 雲端運算, 作業系統, 計算機網路, 資料庫系統, 資訊安全, 軟體工程}

\cventry{Sep 2010–Jun 2014}
        {Bachelor, 資訊工程學系, Score: 3.9}
        {國立成功大學}
        {\url{https://www.ncku.edu.tw/}}
        {}
        {\begin{itemize}
\item 奠定程式設計、演算法與系統基礎
\item 參與校內資訊服務隊，協助維護校園網路與伺服器
\end{itemize}
\textbf{Courses}: 資料結構, 演算法, 作業系統, 計算機組織, 資料庫系統, 計算機網路}
\section{Work}

\cventry{Mar 2021–Present}
        {資深 DevOps 工程師}
        {台灣雲端科技股份有限公司}
        {\url{https://www.twcloudtech.com}}
        {}
        {\begin{itemize}
\item 主導跨 AWS 與 GCP 的多叢集 Kubernetes 平台建置，服務超過 200 個微服務
\item 導入 GitOps（Argo CD）與漸進式交付（Argo Rollouts），部署頻率提升 3 倍
\item 以 Terraform 與 Terragrunt 管理 500+ 雲端資源，實現 100\% 基礎架構即程式碼
\item 建立 Prometheus、Grafana 與 Loki 可觀測性堆疊，將 MTTR 從 45 分鐘降至 20 分鐘
\item 推動 DevSecOps，整合 Trivy、SonarQube 與 OPA 至 CI/CD 流程，降低 60\% 安全風險
\item 指導 5 位初階工程師，建立內部技術文件與 on-call 培訓制度
\end{itemize}
\textbf{Keywords}: Kubernetes, Terraform, GitOps, Argo CD, SRE, AWS, GCP, DevSecOps}

\cventry{Jul 2018–Feb 2021}
        {DevOps 工程師}
        {智聯資訊服務有限公司}
        {\url{https://www.smartlink-info.com.tw}}
        {}
        {\begin{itemize}
\item 建置與維護 Jenkins、GitLab CI 流水線，支援 50+ 專案自動化測試與部署
\item 使用 Docker 與 Helm 將傳統應用容器化，部署時間縮短 70\%
\item 管理 AWS EC2、RDS、S3 與 CloudFront，設計高可用與備援架構
\item 導入 Ansible 進行組態管理，標準化 300+ 台伺服器設定
\item 建立 ELK 日誌分析平台，協助開發團隊快速定位線上問題
\end{itemize}
\textbf{Keywords}: Jenkins, GitLab CI, Docker, Helm, Ansible, AWS, ELK}

\cventry{Aug 2016–Jun 2018}
        {系統維運工程師}
        {宏遠軟體開發股份有限公司}
        {\url{https://www.hongyuan-soft.com}}
        {}
        {\begin{itemize}
\item 負責 Linux 伺服器維運、效能調校與故障排除，維持 99.95\% 服務可用性
\item 撰寫 Shell 與 Python 自動化腳本，減少 50\% 例行維運工時
\item 參與機房搬遷與網路架構優化，確保核心系統零中斷
\item 建置 Nagios 與 Cacti 監控，提前發現潛在容量與效能問題
\end{itemize}
\textbf{Keywords}: Linux, Shell, Python, Nagios, 效能調校, 網路維運}
\section{Languages}

\cvline{Chinese}{Native or Bilingual Proficiency \hfill \textbf{Keywords}: 母語, 技術文件撰寫}
\cvline{English}{Full Professional Proficiency \hfill \textbf{Keywords}: TOEIC 950, 跨}
\section{Skills}

\cvline{雲端平台}{Expert \hfill \textbf{Keywords}: AWS, GCP, Azure, Kubernetes, Terraform}
\cvline{CI/CD 與自動化}{Advanced \hfill \textbf{Keywords}: Jenkins, GitLab CI, Argo CD, Ansible, Helm}
\cvline{監控與可觀測性}{Advanced \hfill \textbf{Keywords}: Prometheus, Grafana, Loki, ELK, Datadog}
\cvline{程式語言與腳本}{Advanced \hfill \textbf{Keywords}: Python, Go, Bash, SQL, YAML}
\cvline{容器與網路}{Expert \hfill \textbf{Keywords}: Docker, Kubernetes, Istio, Cilium, Envoy}
\cvline{資料庫與訊息佇列}{Intermediate \hfill \textbf{Keywords}: PostgreSQL, MySQL, Redis, Kafka, RabbitMQ}
\end{document}